% Section 8 — Related work
\section{Related work}
\label{sec:related}

\paragraph{Learning from outcomes, and paying for search} The lineage in \cref{fig:timeline} ends in
an argument that the next generation of systems should learn mainly from the consequences of their
own actions rather than from human data \citep{silver2025era}, and in the older observation that
general methods which exploit computation, search and learning eventually win
\citep{sutton2019bitter}. AlphaGo Zero is the cleanest demonstration of the half of that argument we
rely on: a network trained on the outcomes of its own search carries much of that search into a
single forward pass \citep{silver2017zero}. The algorithm-discovery systems that followed ---
AlphaTensor, AlphaDev, FunSearch and AlphaEvolve
\citep{fawzi2022alphatensor,mankowitz2023alphadev,romeraparedes2024funsearch,novikov2025alphaevolve}
--- and AutoML-Zero, which discovers complete learning procedures from basic mathematical operations
\citep{real2020automlzero}, all pay for their search when the problem is posed. They are precedents
for outcome-driven discovery, not evidence that every invention needs the same amount of it, and
language models are participants in that lineage, not outsiders to it.

\paragraph{Program synthesis guided by learning} DeepCoder learns properties of programs from
input--output examples to steer an enumerative search \citep{balog2017deepcoder}; DreamCoder grows
a library of abstractions and learns to search with it \citep{ellis2021dreamcoder}; Latent Programmer
first predicts a discrete latent code and then decodes a program from it, a two-level search
\citep{hong2021latent} that is the closest antecedent of our solution embedding. Code-generating
language models sample many programs and filter them on tests \citep{chen2021codex,li2022alphacode}.
In each case learning makes search cheaper; none commits one program before execution as a
contract that a supervisor can audit.

\paragraph{Generating a model from data} Several systems produce a predictor for a dataset in one
pass. HyperNetworks generate the weights of another network \citep{ha2017hypernetworks}; HyperFast
meta-trains a hypernetwork that writes a dataset-specific classifier in a forward pass
\citep{bonet2024hyperfast}; TabPFN predicts by in-context learning, conditioning a pretrained
transformer on the training rows \citep{hollmann2025tabpfn}; and end-to-end transformers predict a
complete mathematical expression, constants included, from observations \citep{kamienny2022e2e}.
These rule out any broad claim of priority for learned data-to-model generation. They differ from
InnovationZero in the object they output: weights, predictions or a formula, rather than a reusable,
inspectable learning procedure with its own \code{fit} and \code{predict} that runs without the
generator.

\paragraph{AutoML and meta-learning} AutoML systems search configuration and pipeline spaces for each
dataset \citep{feurer2015autosklearn,olson2016tpot,wang2021flaml}, and some warm-start that search
from configurations that worked on similar datasets. Warm-starting is amortized proposal over a
closed space, the setting of \cref{prop:closed}: it can choose well, and it cannot propose outside
its portfolio. The engine of \cref{sec:evidence} belongs to this family as a track-S proposer.

\paragraph{The limits of language-model reasoning} A growing literature questions what next-token
training can do on its own: that teacher-forced prediction can fail to learn planning
\citep{bachmann2024pitfalls}, that language models need external verifiers to plan reliably
\citep{kambhampati2024modulo}, and that training and evaluation reward confident guessing
\citep{kalai2025hallucinate}. Our challenge is a direct, if narrow, test of the opposite case:
whether reasoning alone, with no verifier and no execution, produces a working algorithm. It does;
what it produces is the consensus design.

\paragraph{The question that remains} Together these works leave a narrower question than ``can
models generate programs from data''. It is whether a proposer trained on measured outcomes can
commit a useful, reusable fitting procedure for a dataset family it has never seen, before any
candidate-specific feedback, and whether that boundary can be audited independently of the later
fit and evaluation. Supplement~\ref{S-sup:context} discusses the precedents and the debate around
them at greater length.
